\chapter{Introduction}
\label{ch:introduction}

% Guideline recommended length: 1-3 pages (bachelor scale); extended
% here to reflect the greater depth expected of a Master's thesis.
% TODO: read this chapter again once Ch. 7 (Results) is finalised -
% the guideline explicitly asks the Introduction and the Summary to
% mirror each other; any claim made here must be answered there.

\section{Motivation}
\label{sec:motivation}

Autonomous mobile robots are, almost without exception, developed in
simulation first. A simulator offers unlimited, repeatable, and safe
experimentation: a planner can be tuned against a thousand simulated
collisions before it is ever allowed to steer a physical vehicle, and
a software defect that would damage hardware in the real world merely
crashes a process in simulation. This development pattern is so
deeply embedded in contemporary robotics practice that entire
open-source ecosystems -- the Robot Operating System (ROS) and its
successor ROS~2 \cite{quigley2009ros,macenski2022ros2}, the Gazebo
simulator \cite{koenig2004gazebo}, and the Nav2 navigation stack
\cite{macenski2020nav2} used throughout this thesis -- are built
around the assumption that a robot's software stack should run,
largely unmodified, against either a simulated or a physical robot.

That assumption is also where sim-to-real transfer earns its
reputation as one of the harder open problems in applied robotics
\cite{zhao2020sim2real}. A simulator models the physics, sensors, and
actuators its authors chose to model, at the fidelity they chose to
model them; a physical robot is subject to every effect nobody chose
to model at all -- static friction, actuator lag, battery voltage
sag, unsynchronised clocks between distributed compute nodes, and
firmware-level buffering nobody documented. The gap between the two
is not a single bug to be fixed once, but a class of problems that
resurfaces every time a new subsystem crosses from simulation onto
hardware.

This thesis is a case study of that gap, built around a single
physical platform: the Quanser QCar, a 1/10th-scale, real-wheel-drive,
Ackermann-steered research vehicle equipped with a 2D LiDAR, an
inertial measurement unit (IMU), and an onboard NVIDIA Jetson TX2
compute module. The QCar's navigation stack -- SLAM-based mapping,
Adaptive Monte Carlo Localization (AMCL), and Model Predictive Path
Integral (MPPI) trajectory control \cite{williams2016mppi,
williams2017mppi} -- was first developed and validated end-to-end in
a Gazebo simulation, then migrated onto the physical vehicle. The
migration surfaced a sequence of concrete, individually diagnosable
problems: a protocol-level incompatibility that ruled out native
ROS~2 communication with the robot's onboard compute, an uncalibrated
and non-linear mapping from commanded velocity to real actuator duty
cycle, and -- ultimately the most consequential finding of this
project -- a hardware abstraction-layer defect that had, for several
weeks, been mistaken for a fundamental physical property of the
vehicle's drivetrain.

\section{Problem Statement}
\label{sec:problem-statement}

At the outset of this project, the QCar's navigation stack existed
only in simulation: a Gazebo model of the vehicle, driven by the same
Nav2 stack intended for the physical robot, could map an environment,
localise itself within it, and reach goal poses reliably. The
\emph{actual state} was therefore ``a navigation stack validated in
simulation, with no confirmed path to the physical robot''. The
\emph{target state} for this thesis was a navigation stack that
reaches goal poses reliably \emph{on the physical QCar}, with any
sim-to-real discrepancy identified, explained, and -- where
practical -- corrected or compensated for rather than left as an
unexplained residual error.

Reaching that target state required answering four sub-questions,
each of which forms one methodology chapter of this thesis:

\begin{enumerate}
  \item How should the simulation-side navigation stack itself be
        structured so that migrating it to hardware later requires
        changing as little of its Nav2/SLAM configuration as
        possible? (Chapter~\ref{ch:sim-stack})
  \item Given that the robot's onboard compute cannot run the same
        ROS~2 distribution as the navigation stack, how should the two
        be connected so that the rest of the stack does not need to
        know the difference? (Chapter~\ref{ch:hardware-bridge})
  \item How should the mapping from a ROS~2 velocity command to a
        physical actuator signal be identified and calibrated, given
        that the vehicle exhibits both an unknown gain and a hard
        static-friction non-linearity? (Chapter~\ref{ch:calibration})
  \item Once bridged and calibrated, does the navigation stack
        actually reach goals reliably on hardware, and what residual
        sim-to-real gap remains? (Chapters~\ref{ch:nav-integration}
        and~\ref{ch:results})
\end{enumerate}

\section{State of the Art}
\label{sec:state-of-the-art}

ROS~2 was designed specifically to address shortcomings of the
original ROS that made it unsuitable for production and
safety-relevant use, chief among them a lack of real-time support and
a single-point-of-failure master node; it replaces ROS's custom TCP
transport with the Data Distribution Service (DDS) publish-subscribe
middleware standard \cite{omg2015dds,macenski2022ros2,
maruyama2016exploring}. Nav2 is the de facto standard ROS~2
navigation stack, providing pluggable costmaps, global planners, and
local trajectory controllers behind a behaviour-tree-orchestrated
action interface \cite{macenski2020nav2}. MPPI control, the local
trajectory controller used in this project, samples large numbers of
candidate control sequences and weights them by a differentiable cost
function evaluated over a rollout horizon, avoiding the need for a
convex or otherwise analytically tractable vehicle model
\cite{williams2016mppi,williams2017mppi} -- a property well suited to
an Ackermann-steered vehicle whose real dynamics (addressed in
Chapter~\ref{ch:calibration}) turned out to be considerably less
clean than the textbook bicycle model in Chapter~\ref{ch:groundwork}.

Sim-to-real transfer itself is an active research area, most visibly
in reinforcement-learning contexts where a policy trained purely in
simulation is deployed with no further real-world tuning
\cite{zhao2020sim2real}. This thesis does not train a
learning-based policy and therefore sits outside that literature's
core concern (closing the "reality gap" for a black-box learned
controller); its navigation stack -- SLAM, AMCL, and MPPI -- is
classically engineered and identically configured in both simulation
and reality. What this thesis instead documents is a narrower but
still underserved question in the same space: for a classically
engineered stack with an explicit vehicle/sensor/actuator model, what
concrete, individually diagnosable failure modes appear when that
same stack meets real hardware, and how are they found?

\section{Methodological Approach}
\label{sec:methodology-overview}

The work reported in this thesis followed an empirical, iterative
methodology rather than a purely analytical one: each sim-to-real
discrepancy was investigated by direct measurement on the physical
robot (odometry logs, encoder-measured velocity, timestamped Nav2
action status, and, in one case, reading the vendor's hardware
abstraction layer source code directly) rather than inferred solely
from documentation or first-principles modelling. This approach was
not a stylistic preference; Chapter~\ref{ch:calibration} documents a
case in which two independently plausible physical explanations for
an observed 2-second control delay -- drivetrain coasting and a
software command-delivery lag -- were both consistent with the
available second-hand evidence, and only a targeted hardware
measurement, prompted by a direct behavioural observation that
contradicted both prevailing theories, resolved which was correct
(Section~\ref{sec:readmode-bug}).

One control-design attempt in this project was tested directly on the
physical vehicle without first being traced through by hand or
validated in simulation, and it produced a genuine, unplanned
hardware oscillation (Section~\ref{sec:active-braking}). That
incident is reported in full, including why it happened and why it
was reverted rather than iterated on further, because it is as
methodologically relevant to a sim-to-real case study as any of the
project's successful fixes: it demonstrates concretely why validating
control logic before it reaches physical actuators is not merely good
practice but a documented, tested necessity for this class of system.

\section{Thesis Structure}
\label{sec:thesis-structure}

Chapter~\ref{ch:groundwork} introduces the technical foundations this
thesis builds on: the ROS~2/DDS communication model, Ackermann
steering kinematics, the Gazebo simulator and its ROS~2 integration,
SLAM via Cartographer \cite{hess2016cartographer}, and the Nav2
stack's AMCL localisation and MPPI control components.
Chapter~\ref{ch:sim-stack} describes the simulation-side development
of the navigation stack, including two rounds of a systematic vehicle
dynamics bug that recurred across sibling packages, and the discovery
and fix of a URDF parser defect in the underlying middleware.
Chapter~\ref{ch:hardware-bridge} presents the design of the TCP/JSON
bridge that connects the ROS~2-based navigation stack to the robot's
incompatible onboard ROS~2 distribution, including the odometry,
LiDAR angle-convention, and steering-trim fixes required to make the
bridged sensor and actuator data trustworthy. Chapter~\ref{ch:calibration}
covers the identification of the vehicle's throttle response,
including the static-friction deadband and the hardware
abstraction-layer buffering defect that had been mistaken for
physical coasting. Chapter~\ref{ch:nav-integration} describes
integrating and tuning the full Nav2 stack -- including a custom MPPI
critic plugin and AMCL parameter changes -- against the physical
robot, culminating in closed-loop goal-reaching tests.
Chapter~\ref{ch:results} presents and quantifies the resulting
sim-to-real gap across the project's diagnostic measurements, and
Chapter~\ref{ch:summary} answers the research questions posed above
and reflects on what generalises beyond this specific platform.
